% Required packages in the main paper:
% \usepackage{amsmath,amssymb,bm}

\section{Method}
\label{sec:method}

\subsection{Formulation}
\label{sec:formulation}

For $N$ calibrated views, the input is a set of LDR images and their
camera-aligned event streams,
\begin{equation}
    \mathcal{X}=\{(\bm I_i,\mathcal{E}_i,\bm K_i)\}_{i=1}^{N},
    \qquad
    \mathcal{E}_i=\{(\bm u_m,t_m,p_m)\}_{m=1}^{M_i},
    \label{eq:input_definition}
\end{equation}
where $\bm u_m$, $t_m$, and $p_m\in\{+1,-1\}$ denote event position,
timestamp, and polarity. The output is
\begin{equation}
    \mathcal{Y}=\{(D_i,\bm N_i,\bm P_i,\bm T_i)\}_{i=1}^{N},
    \label{eq:output_definition}
\end{equation}
where $D_i$, $\bm N_i$, $\bm P_i$, and $\bm T_i$ are depth, surface
normal, point map, and camera pose, respectively.

The event stream is converted into a polarity-preserving voxel grid with
$B$ temporal bins. With normalized time
$\bar t_m=(B-1)(t_m-t_0)/(t_1-t_0)$, linear temporal splatting gives
\begin{equation}
    V_{i,b}^{p}(\bm u)=
    \sum_{m=1}^{M_i}
    \mathbf{1}[\bm u_m=\bm u]\,
    \mathbf{1}[p_m=p]\,
    \max\!\left(0,1-|b-\bar t_m|\right),
    \quad b\in\{0,\ldots,B-1\}.
    \label{eq:event_voxel_grid}
\end{equation}
We write
\begin{equation}
    \bm V_i=
    [V_{i,0}^{+},\ldots,V_{i,B-1}^{+},
     V_{i,0}^{-},\ldots,V_{i,B-1}^{-}]
    \in\mathbb{R}^{2B\times H\times W}.
    \label{eq:event_voxel_tensor}
\end{equation}

An RGB geometry backbone produces multi-view tokens and coarse geometry,
\begin{equation}
    \bm Z_i,\,D_i^{c},\,\bm P_i^{c},\,\bm T_i
    =\mathcal{G}_{\mathrm{rgb}}
    (\{\bm I_j\}_{j=1}^{N})_i,
    \qquad
    \bm N_i^{c}=\Pi(D_i^{c},\bm K_i),
    \label{eq:rgb_coarse_geometry}
\end{equation}
where $\Pi$ is the differentiable depth-to-normal operator. The complete
mapping is
\begin{equation}
    \boxed{
    \begin{aligned}
        C_i & =\mathcal{R}_{\theta}
        (\bm V_i,\bm I_i,\bm Z_i,D_i^{c},\bm N_i^{c}),\\
        \bm F_i^{e} & =\mathcal{E}_{\phi}(C_i\odot\bm V_i),\\
        \widetilde{\bm N}_i & =
        \mathcal{H}_{N}(\bm F_i^{e},D_i^{c},\bm N_i^{c},C_i),\\
        D_i & =\mathcal{H}_{D}
        (D_i^{c},\bm N_i^{c},\widetilde{\bm N}_i,
         \bm F_i^{e},C_i,\bm K_i),\\
        \bm N_i & =\Pi(D_i,\bm K_i),\\
        \bm P_i & =\mathcal{H}_{P}(\bm Z_i,\bm F_i^{e},C_i).
    \end{aligned}}
    \label{eq:unified_pipeline}
\end{equation}


\subsection{Geometry-Aware Event Reliability}
\label{sec:event_reliability}

During training, controlled event generation provides a complete event stream
$\bm V_i^{\mathrm{full}}$ and its geometry-related component
$\bm V_i^{\mathrm{geo}}$. Their per-pixel temporal masses are
\begin{equation}
    A_i^{\mathrm{full}}(\bm u)=
    \sum_{c=1}^{2B}|V_{i,c}^{\mathrm{full}}(\bm u)|,
    \qquad
    A_i^{\mathrm{geo}}(\bm u)=
    \sum_{c=1}^{2B}|V_{i,c}^{\mathrm{geo}}(\bm u)|.
    \label{eq:event_masses}
\end{equation}
The geometry contribution target is
\begin{equation}
    C_i^{*}(\bm u)=
    M_i(\bm u)\,
    \operatorname{clip}\!\left(
    \frac{A_i^{\mathrm{geo}}(\bm u)}
    {A_i^{\mathrm{full}}(\bm u)+\epsilon},0,1\right),
    \label{eq:contribution_target}
\end{equation}
where $M_i$ is the valid-geometry mask. The contribution network predicts
\begin{equation}
    C_i^{e}=\sigma\!\left(
    \mathcal{R}_{\theta}
    (\bm V_i^{\mathrm{full}},\bm I_i^{e},
     \bm Z_i^{e},D_i^{c,e},\bm N_i^{c,e})\right),
    \quad C_i^{e}\in[0,1]^{H\times W}.
    \label{eq:contribution_prediction}
\end{equation}

For two LDR observations $a$ and $b$ of the same view, the camera, scene
geometry, and event interval are shared. Their contribution predictions are
coupled with the corresponding RGB geometry predictions. We denote the
camera- and scale-aligned geometry tuple by
\begin{equation}
    \bm G_i^{e}=
    [\log\bar D_i^{c,e},\bm N_i^{c,e},\bar{\bm P}_i^{c,e}],
    \qquad e\in\{a,b\},
    \label{eq:paired_geometry_tuple}
\end{equation}
where the bar denotes alignment to the shared camera frame and depth scale.
The per-pixel geometry discrepancy is
\begin{equation}
    d_i^{ab}(\bm u)=
    \lambda_d\rho(\log\bar D_i^{c,a}-
    \operatorname{sg}[\log\bar D_i^{c,b}])+
    \lambda_n\left(1-\left\langle
    \bm N_i^{c,a},\operatorname{sg}[\bm N_i^{c,b}]
    \right\rangle\right)+
    \lambda_p\rho(\bar{\bm P}_i^{c,a}-
    \operatorname{sg}[\bar{\bm P}_i^{c,b}]),
    \label{eq:paired_geometry_discrepancy}
\end{equation}
where all fields on the right-hand side are evaluated at $\bm u$. The
multi-LDR consistency objective is
\begin{equation}
    \mathcal{L}_{\mathrm{LDR}}=
    \frac{
    \sum_{i,\bm u}W_i(\bm u)
    \left|C_i^{a}(\bm u)-
    \operatorname{sg}[C_i^{b}(\bm u)]\right|}
    {\sum_{i,\bm u}W_i(\bm u)+\epsilon}
    +
    \frac{\sum_{i,\bm u}W_i(\bm u)\bar C_i(\bm u)
    d_i^{ab}(\bm u)}
    {\sum_{i,\bm u}W_i(\bm u)\bar C_i(\bm u)+\epsilon},
    \qquad
    \bar C_i=\operatorname{sg}\!\left[\frac{C_i^a+C_i^b}{2}\right],
    \quad W_i=M_i\mathbf{1}[A_i^{\mathrm{full}}>0],
    \label{eq:multi_ldr_consistency}
\end{equation}
where $\operatorname{sg}$ denotes stop-gradient. The controlled-event
regression term is
\begin{equation}
    \mathcal{L}_{\mathrm{dec}}=
    \frac{1}{2}\sum_{e\in\{a,b\}}
    \frac{\sum_{i,\bm u}W_i(\bm u)
    \rho\!\left(C_i^{e}(\bm u)-C_i^{*}(\bm u)\right)}
    {\sum_{i,\bm u}W_i(\bm u)+\epsilon},
    \label{eq:decomposition_supervision}
\end{equation}
where $\rho$ is the smooth-$\ell_1$ penalty.

Let
\begin{equation}
    \Omega_i^{+}=\{\bm u:C_i^{*}(\bm u)\geq\tau_{+}\},
    \qquad
    \Omega_i^{-}=\{\bm u:C_i^{*}(\bm u)\leq\tau_{-}\}.
\end{equation}
The contribution ordering is optimized by
\begin{equation}
    \mathcal{L}_{\mathrm{rank}}=
    \frac{1}{2N}\sum_{i=1}^{N}\sum_{e\in\{a,b\}}
    \left[
    m-mu_{\Omega_i^{+}}(C_i^{e})
    +\mu_{\Omega_i^{-}}(C_i^{e})
    \right]_{+}.
    \label{eq:contribution_ranking}
\end{equation}
The reliability objective is
\begin{equation}
    \boxed{
    \mathcal{L}_{C}=
    \lambda_{\mathrm{dec}}\mathcal{L}_{\mathrm{dec}}+
    \lambda_{\mathrm{LDR}}\mathcal{L}_{\mathrm{LDR}}+
    \lambda_{\mathrm{rank}}\mathcal{L}_{\mathrm{rank}}.}
    \label{eq:reliability_objective}
\end{equation}
At inference, only $C_i=C_i^{e}$ from the available LDR observation is used,
and the selected event tensor is
\begin{equation}
    \widetilde{\bm V}_i=C_i\odot\bm V_i^{\mathrm{full}}.
    \label{eq:reliability_weighted_event}
\end{equation}


\subsection{Event Geometric Property}
\label{sec:event_normal_property}

The weighted voxel grid is encoded with its temporal-bin and polarity axes
preserved:
\begin{equation}
    \bm F_i^{e}=\mathcal{E}_{\phi}(\widetilde{\bm V}_i).
    \label{eq:event_feature_encoding}
\end{equation}
The event branch predicts a bounded local normal residual and a normal
confidence map,
\begin{equation}
    \Delta\bm N_i^{e}=	anh\!\left(
    \mathcal{H}_{\Delta N}
    ([\bm F_i^{e},\log D_i^{c},\bm N_i^{c}])\right),
    \label{eq:event_normal_residual}
\end{equation}
\begin{equation}
    Q_i^{N}=C_i\odot\sigma\!\left(
    \mathcal{H}_{Q}
    ([\bm F_i^{e},\log D_i^{c},\bm N_i^{c}])\right).
    \label{eq:normal_confidence}
\end{equation}
The event-conditioned normal is
\begin{equation}
    \boxed{
    \widetilde{\bm N}_i=
    \operatorname{norm}\!\left(
    \bm N_i^{c}+Q_i^{N}\odot\Delta\bm N_i^{e}
    \right),}
    \label{eq:event_conditioned_normal}
\end{equation}
where $\operatorname{norm}(\bm x)=\bm x/
(\|\bm x\|_2+\epsilon)$.

The event-support mask used for event-normal supervision is
\begin{equation}
    S_i(\bm u)=M_i(\bm u)\,
    \mathbf{1}[A_i^{\mathrm{full}}(\bm u)>0].
    \label{eq:event_normal_support}
\end{equation}
Given $\bm N_i^{*}=\Pi(D_i^{*},\bm K_i)$, the normal-direction term is
\begin{equation}
    \mathcal{L}_{N}^{e}=
    \frac{\sum_{i,\bm u}S_i(\bm u)
    \left(1-langle
    \widetilde{\bm N}_i(\bm u),
    \bm N_i^{*}(\bm u)\rangle\right)}
    {\sum_{i,\bm u}S_i(\bm u)+\epsilon}.
    \label{eq:event_normal_direction_loss}
\end{equation}
The local normal-variation term is
\begin{equation}
    \mathcal{L}_{\nabla N}^{e}=
    \frac{\sum_{i,\bm u}S_i^{\nabla}(\bm u)
    \sum_{d\in\{x,y\}}
    \left\|
    \nabla_d\widetilde{\bm N}_i(\bm u)-
    \nabla_d\bm N_i^{*}(\bm u)
    \right\|_1}
    {\sum_{i,\bm u}S_i^{\nabla}(\bm u)+\epsilon},
    \label{eq:event_normal_derivative_loss}
\end{equation}
where
\begin{equation}
    S_i^{\nabla}(\bm u)=
    S_i(\bm u)S_i(\bm u+\bm e_x)S_i(\bm u+\bm e_y).
    \label{eq:normal_derivative_support}
\end{equation}
Thus, pixels and finite differences without event observations do not
contribute to $\mathcal{L}_{N}^{e}$ or
$\mathcal{L}_{\nabla N}^{e}$.


\subsection{Normal-Guided Geometry Refinement}
\label{sec:normal_guided_refinement}

Depth is initialized by the RGB prediction,
\begin{equation}
    D_i^{(0)}=D_i^{c}.
    \label{eq:depth_initialization}
\end{equation}
At refinement step $k$, its current normal is
\begin{equation}
    \bm N_i^{(k)}=\Pi(D_i^{(k)},\bm K_i).
    \label{eq:iterative_current_normal}
\end{equation}
The shared refinement module predicts a bounded log-depth update,
\begin{equation}
    \delta_i^{(k)}=
    s\tanh\!\left(
    \frac{1}{s}\mathcal{H}_{D}
    ([\bm F_i^{e},\log D_i^{(k)},\log D_i^{c},
      \bm N_i^{(k)},\widetilde{\bm N}_i,Q_i^{N}])
    \right),
    \label{eq:iterative_depth_update}
\end{equation}
followed by
\begin{equation}
    \ell_i^{(k+1)}=
    \operatorname{clip}\!\left(
    \log D_i^{(k)}+\delta_i^{(k)}-\log D_i^{c},
    \log(1-r),\log(1+r)\right),
    \label{eq:bounded_log_ratio}
\end{equation}
\begin{equation}
    D_i^{(k+1)}=D_i^{c}\odot\exp(\ell_i^{(k+1)}),
    \qquad k=0,\ldots,K-1.
    \label{eq:depth_refinement_iteration}
\end{equation}
The final depth and normal are
\begin{equation}
    D_i=D_i^{(K)},
    \qquad
    \bm N_i=\Pi(D_i,\bm K_i).
    \label{eq:final_depth_normal}
\end{equation}
The point-map branch uses the same reliability-weighted event features,
\begin{equation}
    \bm P_i=
    \mathcal{H}_{P}([\bm Z_i,\bm F_i^{e},C_i]).
    \label{eq:point_map_refinement}
\end{equation}

Depth--normal coupling is imposed by
\begin{equation}
    \mathcal{L}_{DN}=
    \frac{\sum_{i,\bm u}M_i(\bm u)Q_i^{N}(\bm u)
    \left(1-\left\langle
    \Pi(D_i,\bm K_i)(\bm u),
    \operatorname{sg}[\widetilde{\bm N}_i(\bm u)]
    \right\rangle\right)}
    {\sum_{i,\bm u}M_i(\bm u)Q_i^{N}(\bm u)+\epsilon}.
    \label{eq:depth_normal_consistency}
\end{equation}
The dense depth update is supervised on the complete valid geometry mask
$M_i$; $S_i$ is used only in the event-normal terms in
Eqs.~\eqref{eq:event_normal_direction_loss}--
\eqref{eq:event_normal_derivative_loss}.


\subsection{Unified Objective}
\label{sec:unified_objective}

The valid finite-difference mask is
\begin{equation}
    M_i^{\nabla}(\bm u)=
    M_i(\bm u)M_i(\bm u+\bm e_x)M_i(\bm u+\bm e_y).
    \label{eq:geometry_stencil_mask}
\end{equation}
The reconstruction terms are
\begin{equation}
    \mathcal{L}_{\mathrm{geo}}=
    \lambda_D\mathcal{L}_{D}+
    \lambda_{\nabla D}\mathcal{L}_{\nabla D}+
    \lambda_{N}\mathcal{L}_{N}+
    \lambda_{P}\mathcal{L}_{P}+
    \lambda_{T}\mathcal{L}_{T},
    \label{eq:geometry_objective}
\end{equation}
with
\begin{align}
    \mathcal{L}_{D}
    &=\frac{\sum_{i,\bm u}M_i(\bm u)
    \rho(D_i(\bm u)-D_i^{*}(\bm u))}
    {\sum_{i,\bm u}M_i(\bm u)+\epsilon},
    \label{eq:depth_loss}\\
    \mathcal{L}_{\nabla D}
    &=\frac{\sum_{i,\bm u}M_i^{\nabla}(\bm u)
    \sum_{d\in\{x,y\}}
    \rho(\nabla_d\log D_i-\nabla_d\log D_i^{*})}
    {\sum_{i,\bm u}M_i^{\nabla}(\bm u)+\epsilon},
    \label{eq:depth_gradient_loss}\\
    \mathcal{L}_{N}
    &=\frac{\sum_{i,\bm u}M_i^{\nabla}(\bm u)
    (1-\langle\bm N_i(\bm u),\bm N_i^{*}(\bm u)\rangle)}
    {\sum_{i,\bm u}M_i^{\nabla}(\bm u)+\epsilon},
    \label{eq:final_normal_loss}\\
    \mathcal{L}_{P}
    &=\frac{\sum_{i,\bm u}M_i(\bm u)
    \|\bm P_i(\bm u)-\bm P_i^{*}(\bm u)\|_1}
    {\sum_{i,\bm u}M_i(\bm u)+\epsilon}.
    \label{eq:point_loss}\\
    \mathcal{L}_{T}
    &=\frac{1}{N}\sum_{i=1}^{N}
    \rho(\bm T_i-\bm T_i^{*}).
    \label{eq:pose_loss}
\end{align}
The complete objective is
\begin{equation}
    \boxed{
    \mathcal{L}=
    \mathcal{L}_{\mathrm{geo}}+
    \mathcal{L}_{C}+
    \lambda_{eN}\mathcal{L}_{N}^{e}+
    \lambda_{e\nabla N}\mathcal{L}_{\nabla N}^{e}+
    \lambda_{DN}\mathcal{L}_{DN}.}
    \label{eq:complete_objective}
\end{equation}

Optimization uses geometry-controlled events to initialize the event geometry
encoder, then learns $C_i$ from full events using
Eq.~\eqref{eq:reliability_objective}, and finally optimizes the
reliability-conditioned normal and geometry refinement using
Eq.~\eqref{eq:complete_objective}. Controlled events, paired LDR observations,
and ground-truth geometry are used only for training. Inference evaluates
Eq.~\eqref{eq:unified_pipeline} from a single LDR multi-view sequence and its
corresponding event streams.
